\documentclass[11pt]{article}

\usepackage[margin=1in]{geometry}
\usepackage[T1]{fontenc}
\usepackage[utf8]{inputenc}
\usepackage{lmodern}
\usepackage{amsmath}
\usepackage{amssymb}
\usepackage{booktabs}
\usepackage{array}
\usepackage{enumitem}
\usepackage{graphicx}
\usepackage{xcolor}
\usepackage{hyperref}

\emergencystretch=2em
\Urlmuskip=0mu plus 2mu

\hypersetup{
  colorlinks=true,
  linkcolor=blue!50!black,
  citecolor=blue!50!black,
  urlcolor=blue!50!black,
  pdftitle={Native-Domain Reservoirs for Mixed Field Reuse},
  pdfauthor={GameCult / Fensalir}
}

\newcommand{\Fensalir}{\textsc{Fensalir}}
\newcommand{\ReSTIR}{ReSTIR}
\newcommand{\AreaReSTIR}{Area ReSTIR}
\newcommand{\TSR}{TSR}
\newcommand{\GRIS}{GRIS}
\newcommand{\Target}{\hat{p}}
\newcommand{\Pdf}{q}
\newcommand{\Measure}{\mu}
\newcommand{\WeightSum}{W}

\title{Native-Domain Reservoirs for Mixed Field Reuse}
\author{GameCult / Fensalir}
\date{Focused Technical Draft, May 2026}

\begin{document}
\maketitle

\begin{abstract}
Reservoir reuse in real-time rendering is usually applied after a visible
surface or pixel sample has already been chosen.  This paper studies a narrower
but important case: mixed field renderers where the reusable sample is the
visible field claim itself.  In \Fensalir, tube fields, SDF surfaces, compact
splat fields, and later density or sensor fields compete for presentation
through one temporal/spatial reservoir.  A texel-owned reservoir can store
field id, depth, color, confidence, and RIS weights, but it cannot prove that a
selected sample remains inside the target domain support after motion,
occlusion, disocclusion, or resolution changes.

We propose a native-domain reservoir contract: every selected sample stores its
producer-domain coordinate, support footprint, proposal measure, target value,
and legal shift metadata.  Pixels are consumers of this selected domain
evidence, not owners of sample identity.  The falsifiable claim is that this
contract reduces temporal leakage and ghosting for mixed tube/SDF/splat fields
compared with texel-owned reuse at similar sample budgets.  We define the
sample contract, target/proposal measures for heterogeneous field claims, reuse
validation rules, required baselines, ablations, and visual diagnostics.  This
draft intentionally does not claim final results; it is the submission-shaped
paper that the implementation must satisfy.
\end{abstract}

\section{Introduction}

\ReSTIR{} and \GRIS{} showed that a renderer can reuse a small number of
representative samples by preserving their target and proposal probability
story \cite{bitterli2020restir,lin2022gris}.  \AreaReSTIR{} sharpened the
lesson for antialiasing and depth of field: when the integration domain is an
area, the selected sample's actual subpixel and lens coordinates are part of
the reservoir state \cite{zhang2024arearestir}.  A reservoir that forgets where
inside the domain its selected sample lives is not allowed to reuse that sample
as though support overlap were obvious.

\Fensalir{} applies that lesson to primary field evidence.  The current engine
contains a shared four-row GPU field reservoir for current-frame RIS, temporal
reuse, spatial/domain reuse, and final resolve.  TubeField and SDF contributors
already share target/source-PDF proposal lanes, and an older
nearest-depth-minus-coverage priority path has been cut.  The live row ABI now
stores selected sample coordinates, producer-domain coordinates, support
footprint, domain kind, and shift kind alongside field id, travel, normal,
guide confidence, proposal weights, and stats.  That makes support overlap a
first-class validation input rather than an inference from a texel.

This paper narrows the previous architecture manifesto into one testable claim:

\begin{quote}
For mixed tube/SDF/splat field rendering, storing selected native-domain
coordinates, support, and legal shift metadata in the reservoir reduces
temporal leakage and ghosting under motion, occlusion, and resolution changes
relative to texel-owned reuse with the same target/source-PDF RIS update.
\end{quote}

The claim is deliberately small enough to kill.  It can fail if the metadata
cost outweighs the quality gain, if support validation rejects too much useful
history, if target measures are incomparable across producers, or if a simpler
\TSR-style resolver plus producer-specific validation performs as well.

\section{Scope}

We evaluate only the core invariant.  The broader \Fensalir{} architecture also
includes semantic grammars, contribution trees, residency scheduling, sensor
fusion, planetary flow, HDR presentation, and D3D12 resource discipline.  Those
systems motivate the domain model but are not the submission claim.

The first evaluated contributors are:

\begin{itemize}[leftmargin=1.5em]
  \item \textbf{TubeField}: texture- or buffer-backed Catmull-Rom curve/tube
  fields with rolling-column identity, curve-local support, emission, and
  explicit motion when available;
  \item \textbf{SDF surface}: bounded analytic or proxy-evaluated signed
  distance surfaces with world/object-domain hit position, normal, material
  class, and travel;
  \item \textbf{compact splat field}: bounded anisotropic Form/Appearance
  splats with local support and projected footprint.
\end{itemize}

Density/extinction, sensor confidence, cube-sphere terrain, and learned
residency are deferred.  They are allowed to influence the data model, not the
result claims.

\section{Related Work}

\paragraph{Reservoir reuse.}
\ReSTIR{} DI and GI provide the practical reservoir shape for real-time
candidate reuse \cite{bitterli2020restir,ouyang2021restirgi}.  \GRIS{}
generalizes the estimator to varied domains and shift mappings
\cite{lin2022gris}.  RTXDI documents production pass boundaries for initial,
temporal, and spatial reuse \cite{rtxdiRestirGI,rtxdiRestirPT}.

\paragraph{Area and reconstruction.}
\AreaReSTIR{} is closest in spirit because it treats subpixel and lens
coordinates as sample identity for area-domain reuse \cite{zhang2024arearestir,
areaRestirCode}.  \TSR{} and modern TAA systems emphasize rejection,
disocclusion classification, sample-age accounting, and debug visualizations
\cite{unrealTSR,unrealTSRFAQ}.  Our distinction is that those diagnostics attach
to selected field evidence before final presentation, not only to color
history.

\paragraph{Field representations.}
Hart's sphere tracing gives the conservative-bound contract for implicit
surfaces \cite{hart1996sphere}.  EWA splatting and 3D Gaussian Splatting show
why anisotropic projected support is a useful primitive
\cite{zwicker2001ewa,ren2002objectewa,kerbl2023gaussian}.  Geometry clipmaps
and Nanite show the value of hierarchy cuts and resident summaries
\cite{asirvatham2005gpuclipmaps,karis2021nanite}.  We do not claim a new
hierarchy algorithm here; we use these systems to constrain the support and
baseline vocabulary.

\section{Native-Domain Reservoir Contract}

A texel-owned field reservoir stores a selected sample as though the pixel were
the native domain.  A native-domain reservoir instead stores a tuple

\begin{equation}
  r = (s, d, u, \Omega_s, \Target(s), \Pdf(s), M, \WeightSum, C, \sigma),
\end{equation}

where \(s\) is the selected sample, \(d\) is its producer domain, \(u\) is the
selected coordinate inside that domain, \(\Omega_s\) is its support footprint,
\(\Target\) is the target importance, \(\Pdf\) is the proposal density or
finite proposal probability, \(M\) is represented candidate count,
\(\WeightSum\) is accumulated weight, \(C\) is final contribution weight, and
\(\sigma\) stores legal shift and validation metadata.

\begin{table}[t]
\centering
\begin{tabular}{p{0.18\linewidth}p{0.27\linewidth}p{0.25\linewidth}p{0.18\linewidth}}
\toprule
Contributor & Native coordinate \(u\) & Support \(\Omega_s\) & Shift metadata \\
\midrule
TubeField & subpixel position, source row/column, curve parameter, segment id
& tube core radius, feather, projected curve support
& rolling-column motion, curve replay/re-eval policy \\
SDF surface & object/world hit position, local object coordinate
& SDF bound radius, normal cone, proxy extent
& object motion, previous hit projection, visibility gate \\
Splat field & local splat center/sample coordinate, orientation frame
& anisotropic envelope, projected footprint
& domain lineage, local-frame transform, Jacobian if measure changes \\
\bottomrule
\end{tabular}
\caption{Minimum native-domain state for the first evaluated contributors.}
\label{tab:sample-contract}
\end{table}

\subsection{Target Measures}

The hard part is not the RIS update; it is defining comparable target measures
for heterogeneous field claims.  We use a projected expected-contribution
target:

\begin{equation}
  \Target_i(c) =
    \int_{\Pi(\Omega_c)}
      \rho_i(x,c)\,
      e_i(x,c)\,
      v_i(x,c)\,
      q_i(x,c)\,dx,
\end{equation}

where \(\Pi(\Omega_c)\) is the projected support of candidate \(c\), \(\rho_i\)
is support/filter density for contributor \(i\), \(e_i\) is estimated visible
field or radiance change, \(v_i\) is visibility/occlusion validity, and
\(q_i\) is producer confidence.  This converts different producer domains into
one unit: expected contribution to the resolved field over projected support.

\begin{table}[t]
\centering
\begin{tabular}{p{0.17\linewidth}p{0.35\linewidth}p{0.34\linewidth}}
\toprule
Contributor & Practical target approximation & Proposal measure \\
\midrule
TubeField & projected core coverage \(\times\) confidence \(\times\)
linear luminance/emission relevance \(\times\) visibility
& finite structural proposal set from proxy/raster coverage or producer curve
samples \\
SDF surface & projected covered area \(\times\) depth/normal confidence
\(\times\) material/radiance relevance \(\times\) visibility
& one deterministic structural proposal per valid hit or finite proxy hit set \\
Splat field & projected anisotropic footprint \(\times\) confidence
\(\times\) Form/Appearance contribution bound
& stochastic or finite selected-cut proposal set over resident splats \\
\bottomrule
\end{tabular}
\caption{Initial target/proposal definitions.  Continuous PDFs are used only
when a producer actually samples from a continuous measure.  Deterministic
structural proposals use explicit finite-set proposal probabilities rather than
a magic free sample.}
\label{tab:target}
\end{table}

For a finite proposal set \(P_i\), the source probability is

\begin{equation}
  \Pdf(c) = \Pr(c \mid c \in P_i),
\end{equation}

usually \(1/|P_i|\) for uniform finite structural proposals, with represented
count preserving how many candidates the selected sample stands for.  A producer
may use \(\Pdf=1\) only for a singleton proposal measure whose represented
count and target are already normalized to that singleton.  This is the part of
the machine that must stay boring and explicit.  Otherwise different producer
families launder incomparable units through the reservoir.

\subsection{Update and Merge}

The update is conventional:

\begin{align}
  w(c) &= \frac{\Target(c)}{\Pdf(c)},\\
  \WeightSum &\leftarrow \WeightSum + w(c),\\
  M &\leftarrow M + m_c,\\
  P(c\text{ selected}) &= \frac{w(c)}{\WeightSum}.
\end{align}

The contribution weight is

\begin{equation}
  C = \frac{\WeightSum}{M\Target(s)}.
\end{equation}

Merging another reservoir \(r_o\) is allowed only after shifting/re-evaluating
its selected sample into the target domain.  The merge probability is

\begin{equation}
  P(r_o\text{ selected}) =
  \frac{\WeightSum_o}{\WeightSum+\WeightSum_o}.
\end{equation}

\section{Reuse Validation}

\subsection{Texel-Owned Baseline}

The baseline stores selected field id, travel/depth, normal, confidence,
target/source-PDF lanes, and RIS stats per pixel row.  Temporal reuse validates
previous row-3 history by reprojected UV, field id, travel, normal, domain
validity, and guide confidence.  Spatial reuse samples neighboring rows and
validates by field id, travel, normal, scalar support, and domain tag.  This is
stronger than naive TAA, but weaker than native-domain reuse because it cannot
ask whether the selected sample coordinate still lies inside target support.

\subsection{Native-Domain Reuse}

Native-domain temporal reuse adds three checks:

\begin{enumerate}[leftmargin=1.5em]
  \item \textbf{Support overlap.}  The previous selected coordinate \(u_p\) is
  transformed into the current domain.  Reuse requires
  \(T_{p\rightarrow c}(u_p) \in \Omega_c\), or a legal shift that creates a
  current-domain coordinate.

  \item \textbf{Producer re-evaluation.}  TubeField and SDF samples are
  re-evaluated at the shifted coordinate.  A stored depth or color is not enough
  authority; the producer-domain sample must still produce a compatible claim.

  \item \textbf{Measure correction.}  If the shift changes measure, the sample
  carries a Jacobian or an explicit finite-set/MIS correction.  If no correction
  is defined, the shift is invalid.
\end{enumerate}

Spatial reuse follows the same rule with neighbor domains: sample first,
validate domain mapping second, merge third.  A neighbor row is never a visible
fallback by itself.

\section{Prototype Evidence}

Figure~\ref{fig:prototype-debug} shows current \Fensalir{} debug captures from
the row-stage reservoir rebuild.  These captures are not the final evaluation:
they show the current diagnostic surface and the failure class that motivates
native-domain state.  The final paper must replace this figure with controlled
side-by-side comparisons and ablations.

\begin{figure}[t]
\centering
\includegraphics[width=0.32\linewidth]{../../artifacts/fensalir-captures/reservoir-rebuild-weighted-final.png}
\includegraphics[width=0.32\linewidth]{../../artifacts/fensalir-captures/reservoir-rebuild-weighted-proposal.png}
\includegraphics[width=0.32\linewidth]{../../artifacts/fensalir-captures/reservoir-rebuild-rejection-codes.png}
\caption{Prototype debug captures from the current shared field reservoir:
final weighted resolve, proposal debug, and rejection codes.  The key
observation is not that these prove the method; they expose why support overlap
and selected producer-domain coordinates must become visible state.}
\label{fig:prototype-debug}
\end{figure}

The first native-domain ABI cut expands each live \texttt{FieldReservoirSample}
row from seven to nine \texttt{float4} lanes:

\begin{center}
\begin{tabular}{ll}
\toprule
Lane & Meaning \\
\midrule
\texttt{colorTravel} & resolved HDR contribution and camera travel \\
\texttt{metadata} & field id and normal/gradient payload \\
\texttt{control} & coverage, scalar support, detail, age/value payload \\
\texttt{guide} & confidence, sample age, domain validity, rejection code \\
\texttt{motion} & previous UV, expected previous travel, explicit-motion flag \\
\texttt{domainSample} & selected pixel/sample UV and packed producer coordinate \\
\texttt{domainSupport} & support footprint, domain kind, shift kind \\
\texttt{proposal} & target, source PDF, represented count, proposal kind \\
\texttt{stats} & selected target, weight sum, candidate count, contribution weight \\
\bottomrule
\end{tabular}
\end{center}

TubeField row-0 proposals write the selected subpixel UV, logical source
column, curve sample coordinate, tube-core support footprint, and explicit
motion shift kind when rolling-column replay is available.  The field id still
uses physical rolling-column identity so age-column history cannot silently
cross-validate.  SDF scene proposals reconstruct the current world hit, pack
object-local coordinates, and mark the sample as object-replayable.  Reservoir
validity now fails when a live sample lacks domain state, and temporal/spatial
validation multiplies the existing field/travel/normal/confidence checks by
support overlap in the selected-domain lane.  The first TubeField replay path
re-evaluates the selected logical column/curve coordinate against the live
producer buffer through a bounded producer-keyed replay manifest and raw source
descriptor table; selected samples choose the matching replay entry by field-id
family instead of relying on a CPU-side single-batch binding.  Manifest overflow
is invalid rather than guessed.
SDF temporal reuse reconstructs the previous selected world hit from stored
sample UV/travel, carries the hit through object center motion, and rejects the
history unless the replayed hit still matches the current camera ray, travel,
and conservative object support.  Exact client SDF distance re-evaluation
remains producer-owned work rather than a generic postprocess authority.
Debug modes expose selected coordinates, support footprint/domain kind, and
shift/rejection state from the same row-3 layer consumed by presentation.
The reservoir is also now a budgeted work grid rather than a native-resolution
buffer: \texttt{GraphicsSettings.FieldReservoirScale} clamps scene evidence,
candidate/history rows, reservoir resolve, bloom sources, and match-window graph
targets below final presentation resolution by default.  The full-resolution
surface is reserved for presentation reconstruction and diagnostics.
The renderer also exposes a field-reservoir mode switch for comparison: native
mode enables domain validity, support overlap, and producer replay during
temporal reuse, while texel-baseline mode keeps the same rows but disables those
native-domain gates.  The first paired capture set
\texttt{reservoir-compare-20260529-231904-*} changed 10.8208\% of final-frame
pixels between modes, and
\texttt{scripts/measure-reservoir-mode-comparison.ps1} now repeats that
measurement.  This proves the comparison switch is live; it is not yet a
temporal leakage or ghosting metric.
The metric harness has since gained a sequence probe: captures may include a
temporal-owner disocclusion debug view derived from previous-UV loss,
previous-closer history, and support-loss rejection codes, and the sequence
measurer uses that mask before falling back to broader rejection output.  This
is still an error-localization surface, not a high-budget reference result.

\section{Required Evaluation}

This paper is not submission-ready until the following evaluation exists.
Reviewers should reject the claim without it.  That sentence stays here until
the figures and tables replace it.

\subsection{Scenes}

\begin{enumerate}[leftmargin=1.5em]
  \item \textbf{Occluded spectrum tubes}: dense rolling TubeField columns
  passing behind/through an SDF occluder.
  \item \textbf{Mixed SDF and splat surface}: an analytic SDF object with
  compact splat detail and camera/subpixel motion.
  \item \textbf{Resolution sweep}: identical field content at multiple output
  resolutions with fixed world-space support.
  \item \textbf{Disocclusion stress}: fast camera/object motion exposing and
  hiding curve/SDF intersections.
  \item \textbf{Near-equal-depth conflict}: overlapping field claims at similar
  travel to test RIS selection versus structural occlusion gates.
\end{enumerate}

\subsection{Baselines}

\begin{itemize}[leftmargin=1.5em]
  \item raw current-frame rendering without temporal reuse;
  \item conventional TAA/TSR-style texel history using depth/normal/field-id
  validation;
  \item texel-owned field reservoir with target/source-PDF lanes but no
  selected native coordinate;
  \item producer-specific TubeField reuse path;
  \item native-domain reservoir with support/shift metadata.
\end{itemize}

\subsection{Ablations}

\begin{itemize}[leftmargin=1.5em]
  \item no selected domain coordinate;
  \item selected coordinate but no support overlap test;
  \item support overlap but no producer re-evaluation;
  \item no measure/Jacobian/MIS correction for shifted samples;
  \item no confidence/variance feedback to reconstruction;
  \item no current-frame structural occlusion gate before RIS.
\end{itemize}

\subsection{Metrics}

The evaluation must report:

\begin{itemize}[leftmargin=1.5em]
  \item temporal leakage pixels in occluded/background regions;
  \item ghosting duration after disocclusion;
  \item support-miss rejection precision/recall against ground-truth producer
  evaluation;
  \item resolved image error against high-sample offline or high-budget
  reference;
  \item stochastic speckle variance over stable regions;
  \item GPU time, memory per reservoir row, bandwidth, and candidate counts;
  \item failure cases where native-domain validation rejects useful history.
\end{itemize}

\section{Novelty Boundary}

The estimator itself is not novel.  The RIS update is standard.  The idea of
shift mappings is not novel after \GRIS.  Area support is not novel after
\AreaReSTIR.  Hierarchical fields are not novel after clipmaps, Nanite, and
splat hierarchies.

The claimed contribution is their boundary in a mixed field renderer:
primary-visible field claims from heterogeneous producers share one reservoir
only if the selected sample stores producer-domain coordinate, support, and
legal shift metadata.  The reservoir is earlier than a lighting/path reservoir
and deeper than a texel history cache.  That is the thing the evaluation must
prove useful.

\section{Limitations}

The contract increases reservoir state.  It may be too expensive for simple
opaque scenes where conventional velocity/depth history already works.  Target
normalization across heterogeneous producers is an open risk; bad targets can
make the reservoir mathematically clean and visually wrong.  Finite structural
proposal measures must be specified carefully, especially when proxy
rasterization creates deterministic candidates.  Finally, native-domain
validation can increase shimmer if it rejects history more aggressively than
the reconstruction pass can repair.

\section{Conclusion}

The paper lives or dies on one invariant: a reusable field reservoir sample
must remember where it was selected in its producer domain and what support it
represents.  Without that, temporal/spatial reuse is a texel-owned cache with
better vocabulary.  With it, mixed tube/SDF/splat fields can be validated,
shifted, rejected, reconstructed, and debugged at the layer where the visible
claim actually lives.  The next implementation pass must build the missing lane
and produce the baselines, ablations, and images listed above.  Until then, this
is a sharp claim with unpaid proof obligations, not a finished paper.

\clearpage
\begin{thebibliography}{20}
\small

\bibitem{bitterli2020restir}
B. Bitterli, C. Wyman, M. Pharr, P. Shirley, A. Lefohn, and W. Jarosz.
\newblock Spatiotemporal reservoir resampling for real-time ray tracing with dynamic direct lighting.
\newblock \emph{ACM Transactions on Graphics}, 2020.
\newblock \url{https://research.nvidia.com/labs/rtr/publication/bitterli2020spatiotemporal/}

\bibitem{ouyang2021restirgi}
Y. Ouyang, S. Liu, M. Kettunen, M. Pharr, and J. Pantaleoni.
\newblock ReSTIR GI: Path resampling for real-time path tracing.
\newblock \emph{Computer Graphics Forum}, 2021.
\newblock \url{https://research.nvidia.com/publication/2021-06_restir-gi-path-resampling-real-time-path-tracing}

\bibitem{lin2022gris}
D. Lin, M. Kettunen, B. Bitterli, J. Pantaleoni, C. Yuksel, and C. Wyman.
\newblock Generalized resampled importance sampling: Foundations of ReSTIR.
\newblock \emph{ACM Transactions on Graphics}, 2022.
\newblock \url{https://graphics.cs.utah.edu/research/projects/gris/sig22_GRIS.pdf}

\bibitem{zhang2024arearestir}
S. Zhang, D. Lin, M. Kettunen, C. Yuksel, and C. Wyman.
\newblock Area ReSTIR: Resampling for real-time defocus and antialiasing.
\newblock \emph{ACM Transactions on Graphics}, 2024.
\newblock \url{https://graphics.cs.utah.edu/research/projects/area-restir/AreaReSTIR-SIGGRAPH2024.pdf}

\bibitem{areaRestirCode}
S. Zhang.
\newblock Area-ReSTIR source code repository.
\newblock \url{https://github.com/guiqi134/Area-ReSTIR}

\bibitem{rtxdiRestirGI}
NVIDIA.
\newblock RTXDI ReSTIR GI documentation.
\newblock \url{https://github.com/NVIDIA-RTX/RTXDI/blob/main/Doc/RestirGI.md}

\bibitem{rtxdiRestirPT}
NVIDIA.
\newblock RTXDI ReSTIR PT documentation.
\newblock \url{https://github.com/NVIDIA-RTX/RTXDI/blob/main/Doc/RestirPT.md}

\bibitem{unrealTSR}
Epic Games.
\newblock Temporal Super Resolution in Unreal Engine.
\newblock \url{https://dev.epicgames.com/documentation/unreal-engine/temporal-super-resolution-in-unreal-engine}

\bibitem{unrealTSRFAQ}
Epic Games.
\newblock Temporal Super Resolution Frequently Asked Questions.
\newblock \url{https://dev.epicgames.com/documentation/unreal-engine/temporal-super-resolution-frequently-asked-questions-for-unreal-engine}

\bibitem{hart1996sphere}
J. C. Hart.
\newblock Sphere tracing: A geometric method for the antialiased ray tracing of implicit surfaces.
\newblock \emph{The Visual Computer}, 1996.
\newblock \url{https://experts.illinois.edu/en/publications/sphere-tracing-a-geometric-method-for-the-antialiased-ray-tracing/}

\bibitem{zwicker2001ewa}
M. Zwicker, H. Pfister, J. van Baar, and M. Gross.
\newblock EWA splatting.
\newblock \emph{IEEE Visualization}, 2001.
\newblock \url{https://dash.harvard.edu/bitstream/handle/1/4138240/Zwicker_EWA.pdf}

\bibitem{ren2002objectewa}
L. Ren, H. Pfister, and M. Zwicker.
\newblock Object space EWA surface splatting.
\newblock \emph{Computer Graphics Forum}, 2002.
\newblock \href{https://publications.ri.cmu.edu/object-space-ewa-surface-splatting-a-hardware-accelerated-approach-to-high-quality-point-rendering/}{CMU publication page}

\bibitem{kerbl2023gaussian}
B. Kerbl, G. Kopanas, T. Leimkuehler, and G. Drettakis.
\newblock 3D Gaussian Splatting for Real-Time Radiance Field Rendering.
\newblock \emph{ACM Transactions on Graphics}, 2023.
\newblock \url{https://arxiv.org/abs/2308.04079}

\bibitem{asirvatham2005gpuclipmaps}
A. Asirvatham and H. Hoppe.
\newblock Terrain rendering using GPU-based geometry clipmaps.
\newblock \emph{GPU Gems 2}, 2005.
\newblock \url{https://developer.nvidia.com/gpugems/gpugems2/part-i-geometric-complexity/chapter-2-terrain-rendering-using-gpu-based-geometry}

\bibitem{karis2021nanite}
B. Karis, R. Stubbe, and G. Wihlidal et al.
\newblock A deep dive into Nanite virtualized geometry.
\newblock \emph{Advances in Real-Time Rendering}, SIGGRAPH 2021.
\newblock \url{https://advances.realtimerendering.com/s2021/Karis_Nanite_SIGGRAPH_Advances_2021_final.pdf}

\end{thebibliography}

\end{document}
